\chapter{KẾT QUẢ THỰC NGHIỆM}

\section{Kết quả thực hiện}

Quá trình nghiên cứu và thực thi trên công cụ AMD Vivado Design Suite đã hoàn thiện thành công hệ thống SoC dựa trên kiến trúc tập lệnh RISC-V và lõi PicoRV32 \cite{riscv_unprivileged,picorv32_repo,amd_vivado_docs}. Về mặt kiến trúc phần cứng, lõi xử lý trung tâm được tích hợp với các mô-đun ngoại vi UART, SPI và I\textsuperscript{2}C thông qua giao diện AXI4-Lite. Hệ thống sử dụng bộ định tuyến Crossbar để điều phối các kênh đọc và ghi độc lập theo mô hình giao dịch của AXI4-Lite \cite{arm_axi4_spec,amd_axi4lite_hls,emil2023axiinterconnect}; các kết quả mô phỏng tích hợp không ghi nhận hiện tượng deadlock trên bus.

Bên cạnh đó, đề tài đã xây dựng hoàn chỉnh khối gia tốc xử lý song song đa luồng, ứng dụng trực tiếp các ô nguyên thủy IDDR và ODDR của FPGA 7 Series để thu và phát dữ liệu ở cả hai cạnh của xung nhịp \cite{amd_ug471}. Hệ thống áp dụng kỹ thuật căn lề trung tâm (Center-aligned Timing), kết hợp với các ràng buộc XDC cho giao tiếp đồng bộ nguồn; cách thiết lập và phân tích các quan hệ định thời này tuân theo hướng dẫn của AMD \cite{amd_ug903,amd_ug949}.

\section{Kết quả thử nghiệm}

\subsection{Môi trường và phương pháp thử nghiệm}
Môi trường thử nghiệm chức năng và tổng hợp vật lý (RTL Implementation) được triển khai trên phần mềm AMD Vivado v2025.2, hướng tới FPGA Artix-7 \path{xc7a200tfbg676-2L} \cite{amd_vivado_docs,amd_ds180,amd_ds181}. Tần số hoạt động của hệ thống được thiết lập ở mức $100\text{ MHz}$ trong mô phỏng chức năng và $50\text{ MHz}$ đối với các báo cáo sau hiện thực. Với cấu hình mô phỏng $100\text{ MHz}$, việc truyền một mẫu ở mỗi cạnh clock tương ứng tốc độ dữ liệu lý thuyết $200\text{ Mbps/pin}$ \cite{jedec_ddr,amd_ug471}. Các tiêu chí đánh giá gồm tính chính xác của dạng sóng, mức sử dụng tài nguyên, công suất ước tính và các biên an toàn định thời.

\subsection{Kiểm chứng bộ điều khiển UART}
Ở kịch bản kiểm tra tốc độ baud và cấu trúc khung truyền, giao diện AXI ghi giá trị $0\text{x}55$ vào thanh ghi \texttt{REG\_TXDATA}. Dạng sóng trên đường \texttt{uart\_tx} xuất hiện đúng Start Bit, tám bit dữ liệu theo thứ tự LSB-first và Stop Bit theo cấu trúc khung UART thông dụng \cite{ti_uart_overview,huang2024uartcommunication}; thời gian mỗi bit phù hợp với tốc độ baud đã cấu hình.

\begin{figure}[H]
\centering
\includegraphics[width=\linewidth,height=0.34\textheight,keepaspectratio]{TB/UART_WAVE1.png}
\caption{Dạng sóng kiểm tra tốc độ baud và cấu trúc khung truyền UART.}
\label{fig:tb-uart-wave1}
\end{figure}

\begin{figure}[H]
\centering
\includegraphics[width=\linewidth,height=0.18\textheight,keepaspectratio]{TB/UART_LOG1.png}
\caption{Kết quả mô phỏng kiểm tra tốc độ baud và cấu trúc khung truyền UART.}
\label{fig:tb-uart-log1}
\end{figure}

Phép thử với mẫu $0\text{xA5}$ cho thấy khung truyền phát hoàn chỉnh và khớp với dữ liệu mong muốn, xác nhận máy trạng thái TX hoạt động ổn định.

\begin{figure}[H]
\centering
\includegraphics[width=\linewidth,height=0.34\textheight,keepaspectratio]{TB/UART_WAVE2.png}
\caption{Dạng sóng truyền mẫu dữ liệu UART TX $0\text{xA5}$.}
\label{fig:tb-uart-wave2}
\end{figure}

\begin{figure}[H]
\centering
\includegraphics[width=\linewidth,height=0.18\textheight,keepaspectratio]{TB/UART_LOG2.png}
\caption{Kết quả mô phỏng truyền mẫu dữ liệu UART TX $0\text{xA5}$.}
\label{fig:tb-uart-log2}
\end{figure}

Thanh ghi \texttt{REG\_STATUS} cũng phản ánh đúng trạng thái truyền: cờ \texttt{Tx Ready} chuyển từ $1$ (IDLE) xuống $0$ (BUSY) khi bắt đầu và trở lại $1$ sau khi khung truyền hoàn tất.

\begin{figure}[H]
\centering
\includegraphics[width=\linewidth,height=0.34\textheight,keepaspectratio]{TB/UART_WAVE3.png}
\caption{Dạng sóng kiểm tra thanh ghi trạng thái Ready/Busy của UART.}
\label{fig:tb-uart-wave3}
\end{figure}

\begin{figure}[H]
\centering
\includegraphics[width=\linewidth,height=0.18\textheight,keepaspectratio]{TB/UART_LOG3.png}
\caption{Kết quả mô phỏng kiểm tra thanh ghi trạng thái Ready/Busy của UART.}
\label{fig:tb-uart-log3}
\end{figure}

Đối với chức năng RX, testbench đưa khung $0\text{x}3\text{C}$ vào \texttt{uart\_rx}. Sau khi thu hoàn tất, chu kỳ đọc AXI tại \texttt{REG\_RXDATA} trả về đúng giá trị $0\text{x}3\text{C}$.

\begin{figure}[H]
\centering
\includegraphics[width=\linewidth,height=0.34\textheight,keepaspectratio]{TB/UART_WAVE4.png}
\caption{Dạng sóng kiểm tra chức năng nhận dữ liệu UART RX $0\text{x3C}$.}
\label{fig:tb-uart-wave4}
\end{figure}

\begin{figure}[H]
\centering
\includegraphics[width=\linewidth,height=0.18\textheight,keepaspectratio]{TB/UART_LOG4.png}
\caption{Kết quả mô phỏng kiểm tra chức năng nhận dữ liệu UART RX $0\text{x3C}$.}
\label{fig:tb-uart-log4}
\end{figure}

\subsection{Kiểm chứng bộ điều khiển SPI}
Trước khi giao dịch bắt đầu, đường \texttt{spi\_sclk} duy trì ở mức logic $0$, xác nhận cấu hình CPOL=0. Trong quá trình truyền song công, mẫu $0\text{xD6}$ được phát trên \texttt{spi\_mosi}, đồng thời bộ điều khiển thu đúng giá trị $0\text{x3C}$ từ \texttt{spi\_miso}; cách tổ chức các đường SCLK, MOSI, MISO và chọn chip phù hợp với mô hình giao tiếp SPI \cite{analog_spi_overview,digilent_spi}.

\begin{figure}[H]
\centering
\includegraphics[width=\linewidth,height=0.34\textheight]{TB/SPI_WAVE1-2.png}
\caption{Dạng sóng kiểm tra CPOL và truyền nhận song công SPI.}
\label{fig:tb-spi-wave1-2}
\end{figure}

\begin{figure}[H]
\centering
\includegraphics[width=\linewidth,height=0.18\textheight,keepaspectratio]{TB/SPI_LOG1-2.png}
\caption{Kết quả mô phỏng kiểm tra CPOL và truyền nhận song công SPI.}
\label{fig:tb-spi-log1-2}
\end{figure}

Bên cạnh tín hiệu truyền nhận, cơ chế chọn chip (Chip Select) cũng được kiểm chứng thông qua thanh ghi \texttt{REG\_SPI\_CS} (Hình~\ref{fig:tb-spi-wave3}). Dạng sóng minh chứng rằng các thao tác ghi giá trị $0$ hoặc $1$ từ bus AXI phản ánh tức thời và chính xác lên chân vật lý \texttt{spi\_cs}, đồng thời kênh đọc AXI cũng trả về đúng giá trị cấu hình, đảm bảo tính linh hoạt trong điều khiển hệ thống đa tớ (multi-slave).

\begin{figure}[H]
\centering
\includegraphics[width=\linewidth,height=0.34\textheight]{TB/SPI_WAVE3.png}
\caption{Dạng sóng kiểm tra thao tác thanh ghi Chip Select của SPI.}
\label{fig:tb-spi-wave3}
\end{figure}

\begin{figure}[H]
\centering
\includegraphics[width=\linewidth,height=0.18\textheight,keepaspectratio]{TB/SPI_LOG3.png}
\caption{Kết quả mô phỏng kiểm tra thao tác thanh ghi Chip Select của SPI.}
\label{fig:tb-spi-log3}
\end{figure}

\subsection{Kiểm chứng bộ điều khiển \texorpdfstring{I\textsuperscript{2}C}{I2C}}
Giao thức I\textsuperscript{2}C sử dụng hai đường SCL và SDA kiểu cực máng hở, đồng thời quy định rõ các điều kiện START, STOP và chu kỳ ACK/NACK \cite{nxp_um10204}. Trong phép thử truy xuất thanh ghi, thao tác ghi và đọc giá trị $0\text{xA7}$ qua AXI hoàn tất chính xác.

\begin{figure}[H]
\centering
\includegraphics[width=\linewidth,height=0.34\textheight,keepaspectratio]{TB/I2C_WAVE1.png}
\caption{Dạng sóng kiểm tra thao tác đọc/ghi thanh ghi I\textsuperscript{2}C qua AXI.}
\label{fig:tb-i2c-wave1}
\end{figure}

\begin{figure}[H]
\centering
\includegraphics[width=\linewidth,height=0.18\textheight,keepaspectratio]{TB/I2C_LOG1.png}
\caption{Kết quả mô phỏng kiểm tra thao tác đọc/ghi thanh ghi I\textsuperscript{2}C qua AXI.}
\label{fig:tb-i2c-log1}
\end{figure}

Về phương diện tạo tín hiệu điều khiển trên bus, mạch FSM đã tự động tạo điều kiện START chuẩn xác bằng cách hạ đường \texttt{sda\_out} xuống mức thấp trong khi \texttt{scl\_out} vẫn duy trì ở mức cao (Hình~\ref{fig:tb-i2c-wave2}). Xuyên suốt quá trình truyền byte dữ liệu $0\text{x}3\text{C}$ (Hình~\ref{fig:tb-i2c-wave3-ack} và Hình~\ref{fig:tb-i2c-wave3-nack}), hệ thống tự động thả nổi bus (chuyển \texttt{sda\_oe} sang trạng thái ngắt) tại chu kỳ xung nhịp thứ 9 để thu thập tín hiệu phản hồi. Khi Slave kéo \texttt{sda\_in} xuống $0$, cờ trạng thái báo nhận thành công tín hiệu ACK; ngược lại, khi \texttt{sda\_in} giữ nguyên mức $1$, hệ thống nhận diện tức thời tín hiệu NACK. 

\begin{figure}[H]
\centering
\includegraphics[width=\linewidth,height=0.34\textheight]{TB/I2C_WAVE2.png}
\caption{Dạng sóng kiểm tra điều kiện START của giao tiếp I\textsuperscript{2}C.}
\label{fig:tb-i2c-wave2}
\end{figure}

\begin{figure}[H]
\centering
\includegraphics[width=\linewidth,height=0.18\textheight,keepaspectratio]{TB/I2C_LOG2.png}
\caption{Kết quả mô phỏng kiểm tra điều kiện START của giao tiếp I\textsuperscript{2}C.}
\label{fig:tb-i2c-log2}
\end{figure}

\begin{figure}[H]
\centering
\includegraphics[width=\linewidth,height=0.34\textheight]{TB/I2C_WAVE3_ACK.png}
\caption{Dạng sóng truyền byte và nhận phản hồi ACK của I\textsuperscript{2}C.}
\label{fig:tb-i2c-wave3-ack}
\end{figure}

\begin{figure}[H]
\centering
\includegraphics[width=\linewidth,height=0.18\textheight,keepaspectratio]{TB/I2C_LOG3_ACK.png}
\caption{Kết quả mô phỏng truyền byte và nhận phản hồi ACK của I\textsuperscript{2}C.}
\label{fig:tb-i2c-log3-ack}
\end{figure}

\begin{figure}[H]
\centering
\includegraphics[width=\linewidth,height=0.34\textheight]{TB/I2C_WAVE3_NACK.png}
\caption{Dạng sóng truyền byte và nhận phản hồi NACK của I\textsuperscript{2}C.}
\label{fig:tb-i2c-wave3-nack}
\end{figure}

\begin{figure}[H]
\centering
\includegraphics[width=\linewidth,height=0.18\textheight,keepaspectratio]{TB/I2C_LOG3_NACK.png}
\caption{Kết quả mô phỏng truyền byte và nhận phản hồi NACK của I\textsuperscript{2}C.}
\label{fig:tb-i2c-log3-nack}
\end{figure}

Cuối cùng, giao dịch I\textsuperscript{2}C được đóng gói an toàn bằng điều kiện STOP (Hình~\ref{fig:tb-i2c-wave4}), trong đó tín hiệu \texttt{sda\_out} nảy lên mức cao trong thời điểm \texttt{scl\_out} giữ trạng thái tích cực, đáp ứng hoàn hảo tiêu chuẩn của giao thức.

\begin{figure}[H]
\centering
\includegraphics[width=\linewidth,height=0.34\textheight]{TB/I2C_WAVE4.png}
\caption{Dạng sóng kiểm tra điều kiện STOP của giao tiếp I\textsuperscript{2}C.}
\label{fig:tb-i2c-wave4}
\end{figure}

\begin{figure}[H]
\centering
\includegraphics[width=\linewidth,height=0.18\textheight,keepaspectratio]{TB/I2C_LOG4.png}
\caption{Kết quả mô phỏng kiểm tra điều kiện STOP của giao tiếp I\textsuperscript{2}C.}
\label{fig:tb-i2c-log4}
\end{figure}

\subsection{Kiểm chứng khối gia tốc xử lý song song DDR}
Khối gia tốc song song ứng dụng kỹ thuật DDR, trong đó IDDR và ODDR lần lượt chốt hoặc phát dữ liệu tại cả cạnh lên và cạnh xuống của clock \cite{amd_ug471,jedec_ddr}. Trong kịch bản kích hoạt đồng bộ (Hình~\ref{fig:tb-ddr-wave1}), ngay khi nhận lệnh điều khiển toàn cục từ CPU, 4 luồng tính toán song song xử lý các toán hạng và tự động xả chuỗi 4 mẫu kết quả qua bộ tuần tự hóa ODDR. Dạng sóng minh họa chuỗi giá trị $0\text{x}11$, $0\text{x}22$, $0\text{x}33$ và $0\text{x}44$ được phát chính xác; cờ \texttt{result\_valid} tự xóa sau khi kênh đọc AXI hoàn tất truy xuất.

\begin{figure}[H]
\centering
\includegraphics[width=\linewidth,height=0.34\textheight]{TB/DDR_WAVE1.png}
\caption{Dạng sóng kiểm tra kích hoạt đồng bộ toàn bộ các làn xử lý DDR.}
\label{fig:tb-ddr-wave1}
\end{figure}

\begin{figure}[H]
\centering
\includegraphics[width=\linewidth,height=0.22\textheight,keepaspectratio]{TB/DDR_LOG1.png}
\caption{Kết quả mô phỏng kích hoạt đồng bộ toàn bộ các làn xử lý DDR.}
\label{fig:tb-ddr-log1}
\end{figure}

Khả năng kích hoạt độc lập từng làn được kiểm tra tại Hình~\ref{fig:tb-ddr-wave2}. Kết quả cho thấy mỗi làn chỉ tiếp nhận và xử lý dữ liệu khi tín hiệu kích hoạt tương ứng được xác lập; các làn còn lại giữ nguyên trạng thái, qua đó xác nhận cơ chế điều khiển độc lập không gây tác động chéo.

\begin{figure}[H]
\centering
\includegraphics[width=\linewidth,height=0.34\textheight]{TB/DDR_WAVE2.png}
\caption{Dạng sóng kiểm tra khả năng kích hoạt độc lập từng làn xử lý DDR.}
\label{fig:tb-ddr-wave2}
\end{figure}

\begin{figure}[H]
\centering
\includegraphics[width=\linewidth,height=0.22\textheight,keepaspectratio]{TB/DDR_LOG2.png}
\caption{Kết quả mô phỏng khả năng kích hoạt độc lập từng làn xử lý DDR.}
\label{fig:tb-ddr-log2}
\end{figure}

Tính toàn vẹn tín hiệu đa sườn xung được khẳng định qua phép thử cách ly sườn xung (Edge Isolation) tại Hình~\ref{fig:tb-ddr-wave3}. Khi hệ thống nạp dữ liệu $0\text{x}999$ vào sườn lên (Rising Edge) và $0\text{x}000$ vào sườn xuống (Falling Edge) thông qua IDDR, mạch phân luồng bóc tách dữ liệu tuyệt đối chính xác. Luồng sườn lên xử lý độc lập mà không gây rò rỉ tín hiệu hay ghi đè nhầm sang luồng sườn xuống, minh chứng cho sự phối hợp tin cậy giữa các tầng thanh ghi trễ đồng bộ (Pipeline Matching).

\begin{figure}[H]
\centering
\includegraphics[width=\linewidth,height=0.34\textheight]{TB/DDR_WAVE3.png}
\caption{Dạng sóng kiểm tra khả năng cách ly dữ liệu giữa hai sườn xung DDR.}
\label{fig:tb-ddr-wave3}
\end{figure}

\begin{figure}[H]
\centering
\includegraphics[width=\linewidth,height=0.22\textheight,keepaspectratio]{TB/DDR_LOG3.png}
\caption{Kết quả mô phỏng khả năng cách ly dữ liệu giữa hai sườn xung DDR.}
\label{fig:tb-ddr-log3}
\end{figure}

\subsubsection{Đánh giá thông lượng thông qua Console Log (Throughput Benchmark)}
Nhằm định lượng mức độ hiệu quả của kiến trúc, một hệ thống giám sát đo lường (Benchmark) đã được tích hợp để so sánh thời gian thực thi tác vụ truyền tải 4 mẫu dữ liệu giữa kiến trúc 1 Lane truyền thống và kiến trúc đa luồng DDR 2 Lane. Dạng sóng đo độ trễ và kết quả thực thi trên cửa sổ Console Log của Vivado Simulator được trình bày lần lượt tại Hình~\ref{fig:tb-ddr-wave4} và Hình~\ref{fig:tb-ddr-log4}.

\begin{figure}[H]
\centering
\includegraphics[width=\linewidth,height=0.34\textheight]{TB/DDR_WAVE4.png}
\caption{Dạng sóng so sánh độ trễ giữa kiến trúc 1 Lane và kiến trúc DDR 2 Lane.}
\label{fig:tb-ddr-wave4}
\end{figure}

\begin{figure}[H]
\centering
\includegraphics[width=\linewidth,height=0.3\textheight,keepaspectratio]{TB/DDR_LOG4.png}
\caption{Kết quả đánh giá thông lượng và hệ số tăng tốc của kiến trúc DDR.}
\label{fig:tb-ddr-log4}
\end{figure}

Kết quả đo cho thấy kiến trúc DDR 2 Lane hoàn thành tác vụ trong $110\text{ ns}$, so với $320\text{ ns}$ của kiến trúc 1 Lane. Hệ số tăng tốc đạt $2.91$ lần và thời gian xử lý giảm $65.6\%$, qua đó xác nhận lợi thế thông lượng của kiến trúc DDR.

\subsection{Kiểm chứng mức tích hợp hệ thống (SoC Integration)}
Môi trường mô phỏng SoC được sử dụng để xác thực khả năng giải mã địa chỉ và vận hành đồng thời nhiều ngoại vi. Kết quả tại Hình~\ref{fig:tb-soc-wave1} và Hình~\ref{fig:tb-soc-log1} cho thấy AXI Crossbar định tuyến đúng giao dịch đến thiết bị mục tiêu mà không gây tranh chấp bus.

\begin{figure}[H]
\centering
\includegraphics[width=\linewidth,height=0.34\textheight]{TB/SOC_WAVE1.png}
\caption{Dạng sóng kiểm tra truy cập và giải mã địa chỉ AXI của hệ thống SoC.}
\label{fig:tb-soc-wave1}
\end{figure}

\begin{figure}[H]
\centering
\includegraphics[width=\linewidth,height=0.18\textheight,keepaspectratio]{TB/SOC_LOG1.png}
\caption{Kết quả mô phỏng kiểm tra truy cập và giải mã địa chỉ AXI của hệ thống SoC.}
\label{fig:tb-soc-log1}
\end{figure}

Các kịch bản tiếp theo kích hoạt đồng thời UART và SPI, sau đó mở rộng sang UART, SPI và I\textsuperscript{2}C. Các dạng sóng tại Hình~\ref{fig:tb-soc-wave2} và Hình~\ref{fig:tb-soc-wave3} cho thấy các ngõ ra vật lý hoạt động đan xen trong cùng một khoảng thời gian mà không làm mất dữ liệu.

\begin{figure}[H]
\centering
\includegraphics[width=\linewidth,height=0.34\textheight]{TB/SOC_WAVE2.png}
\caption{Dạng sóng kiểm tra UART và SPI hoạt động đồng thời.}
\label{fig:tb-soc-wave2}
\end{figure}

\begin{figure}[H]
\centering
\includegraphics[width=\linewidth,height=0.18\textheight,keepaspectratio]{TB/SOC_LOG2.png}
\caption{Kết quả mô phỏng kiểm tra UART và SPI hoạt động đồng thời.}
\label{fig:tb-soc-log2}
\end{figure}

\begin{figure}[H]
\centering
\includegraphics[width=\linewidth,height=0.34\textheight]{TB/SOC_WAVE3.png}
\caption{Dạng sóng kiểm tra UART, SPI và I\textsuperscript{2}C hoạt động đồng thời.}
\label{fig:tb-soc-wave3}
\end{figure}

\begin{figure}[H]
\centering
\includegraphics[width=\linewidth,height=0.18\textheight,keepaspectratio]{TB/SOC_LOG3.png}
\caption{Kết quả mô phỏng kiểm tra UART, SPI và I\textsuperscript{2}C hoạt động đồng thời.}
\label{fig:tb-soc-log3}
\end{figure}

Trong kịch bản giao dịch đan xen, chuỗi lệnh đọc và ghi thay đổi liên tục giữa các vùng địa chỉ UART, SPI và I\textsuperscript{2}C. Kết quả tại Hình~\ref{fig:tb-soc-wave4} và Hình~\ref{fig:tb-soc-log4} xác nhận cơ chế chốt mục tiêu của AXI Crossbar duy trì đúng đường phản hồi và không xảy ra deadlock.

\begin{figure}[H]
\centering
\includegraphics[width=\linewidth,height=0.34\textheight]{TB/SOC_WAVE4.png}
\caption{Dạng sóng kiểm tra các giao dịch AXI đan xen giữa nhiều ngoại vi.}
\label{fig:tb-soc-wave4}
\end{figure}

\begin{figure}[H]
\centering
\includegraphics[width=\linewidth,height=0.26\textheight,keepaspectratio]{TB/SOC_LOG4.png}
\caption{Kết quả mô phỏng kiểm tra các giao dịch AXI đan xen giữa nhiều ngoại vi.}
\label{fig:tb-soc-log4}
\end{figure}

\subsection{Kết quả tổng hợp và thực thi phần cứng (RTL Implementation)}

\subsubsection{Phân tích tài nguyên phần cứng}
Báo cáo sau định tuyến của Vivado cho thấy hệ thống sử dụng $447$ Slice LUTs, $573$ Slice Registers và $213$ Slice. Khối phát DDR sử dụng $33$ ô OLOGIC tại các chân I/O; đây là nhóm tài nguyên chuyên dụng chứa các phần tử tuần tự ngõ ra của FPGA 7 Series \cite{amd_ug471}.

\begin{table}[htbp]
\centering
\setlength{\tabcolsep}{1pt}
\caption{Bảng tổng hợp sử dụng tài nguyên phần cứng thực tế trên FPGA Artix-7 (xc7a200tfbg676-2L).}
\label{tab:resource-utilization}
\begin{tabular}{|>{\raggedright\arraybackslash}p{5.5cm}|>{\centering\arraybackslash}p{2.2cm}|>{\centering\arraybackslash}p{2.5cm}|>{\centering\arraybackslash}p{1.8cm}|}
\hline
\textbf{Tài nguyên (Resource)} & \textbf{Sử dụng (Used)} & \textbf{Khả dụng (Available)} & \textbf{Tỷ lệ (\%)} \\ \hline
Slice & 213 & 33,650 & 0.63\% \\ \hline
Slice LUTs (Total) & 447 & 134,600 & 0.33\% \\ \hline
--- \textit{LUT as Logic} & 415 & 134,600 & 0.31\% \\ \hline
--- \textit{LUT as Distributed RAM} & 32 & 46,200 & 0.07\% \\ \hline
Slice Registers (FFs) & 573 & 269,200 & 0.21\% \\ \hline
Block RAM Tile (RAMB36E1) & 0 & 365 & 0.00\% \\ \hline
Bonded IOB & 78 & 400 & 19.50\% \\ \hline
--- \textit{OLOGIC (ODDR / Output FF)} & 33 & 400 & 8.25\% \\ \hline
Global Clock Buffer (BUFGCTRL) & 1 & 32 & 3.13\% \\ \hline
\end{tabular}
\end{table}

Thiết kế chỉ chiếm $0.33\%$ tổng số Slice LUTs, $0.21\%$ thanh ghi và $0.63\%$ số Slice của FPGA mục tiêu. Các giá trị khả dụng dùng để tính tỷ lệ được đối chiếu với thông tin tài nguyên của họ FPGA 7 Series và Artix-7 \cite{amd_ds180,amd_ds181}. Thiết kế không sử dụng DSP và Block RAM Tile, do đó các tài nguyên này vẫn sẵn sàng cho những khối mở rộng.

\subsubsection{Phân tích định thời và công suất tiêu thụ}
Báo cáo định thời tĩnh tại tần số $50\text{ MHz}$ ghi nhận $WNS=+12.339\text{ ns}$ và $WHS=+0.105\text{ ns}$. Cả hai giá trị đều dương và không có endpoint vi phạm, cho thấy thiết kế đạt timing closure ở cấu hình khảo sát; cách diễn giải các chỉ số slack và quy trình kiểm tra timing tuân theo phương pháp luận của Vivado \cite{amd_ug903,amd_ug949}.

\begin{table}[htbp]
\centering
\caption{Bảng tổng hợp kết quả phân tích định thời của hệ thống SoC tại tần số $50\text{ MHz}$.}
\label{tab:timing-summary}
\begin{tabular}{|>{\raggedright\arraybackslash}p{6cm}|>{\centering\arraybackslash}p{3cm}|>{\centering\arraybackslash}p{3cm}|}
\hline
\textbf{Chỉ số định thời} & \textbf{Giá trị (ns)} & \textbf{Trạng thái} \\ \hline
Worst Negative Slack (WNS) & $+12.339$ & Đạt \\ \hline
Worst Hold Slack (WHS) & $+0.105$ & Đạt \\ \hline
Số endpoint vi phạm & 0 & Đạt \\ \hline
\end{tabular}
\end{table}

\begin{table}[htbp]
\centering
\setlength{\tabcolsep}{1pt}
\caption{Bảng tổng hợp công suất tiêu thụ hệ thống SoC tại tần số $50\text{ MHz}$.}
\label{tab:power-summary}
\begin{tabular}{|>{\raggedright\arraybackslash}p{4.8cm}|>{\centering\arraybackslash}p{2.2cm}|>{\centering\arraybackslash}p{2.2cm}|>{\centering\arraybackslash}p{2.5cm}|}
\hline
\textbf{Thành phần công suất} & \textbf{Công suất (W)} & \textbf{Tỷ lệ tổng (\%)} & \textbf{Tỷ lệ động (\%)} \\ \hline
\textbf{Công suất động} & 0.051 & 28.02 & 100 \\ \hline
--- \textit{Clock} & 0.002 & 1.10 & 4 \\ \hline
--- \textit{Signals} & $<0.001$ & $<0.55$ & $<2$ \\ \hline
--- \textit{Logic} & $<0.001$ & $<0.55$ & $<2$ \\ \hline
--- \textit{BRAM} & 0 & 0.00 & 0 \\ \hline
--- \textit{I/O} & 0.047 & 25.82 & 92 \\ \hline
Công suất tĩnh (Device Static) & 0.131 & 71.98 & --- \\ \hline
Tổng công suất trên chip & 0.182 & 100.00 & --- \\ \hline
\end{tabular}
\end{table}

Theo báo cáo ước tính công suất của Vivado, tổng công suất trên chip là $0.182\text{ W}$, gồm $0.051\text{ W}$ công suất động và $0.131\text{ W}$ công suất tĩnh \cite{amd_vivado_docs}. Thành phần I/O tiêu thụ $0.047\text{ W}$, tương ứng khoảng $92\%$ công suất động, và là thành phần động lớn nhất của thiết kế.

\section{Đánh giá và thảo luận}
Thông qua các dữ liệu thực nghiệm trên, hệ thống SoC đạt sự cân bằng giữa hiệu suất, công suất và diện tích. Kiến trúc AXI4-Lite cung cấp giao diện ánh xạ thanh ghi với các kênh đọc/ghi tách biệt, phù hợp để điều phối nhiều mô-đun ngoại vi \cite{arm_axi4_spec,emil2023axiinterconnect}. Khối gia tốc DDR đa luồng đạt hệ số tăng tốc $2.91$ lần trong kịch bản benchmark của đề tài; kết quả này phù hợp với nguyên lý tăng số lần truyền dữ liệu trên mỗi chu kỳ của liên kết DDR \cite{psarras2017nocddr,jedec_ddr}.

Phương pháp căn lề trung tâm tại các chân phát dữ liệu DDR cho thấy tính thực tiễn trong cấu hình khảo sát. Việc đặt cạnh lấy mẫu vào vùng dữ liệu ổn định giúp cải thiện biên setup/hold; tuy nhiên, quan hệ pha, độ trễ đường truyền và các ràng buộc XDC vẫn phải được mô tả đầy đủ cho từng giao diện cụ thể \cite{arumugham2000ddrtiming,amd_ug903,amd_ug471}.

\section{Hạn chế của đề tài}
Mặc dù hệ thống đáp ứng các chỉ tiêu ở cấu hình khảo sát, một số khía cạnh vẫn cần được nâng cấp. Thứ nhất, khối cộng tổ hợp trong mô-đun xử lý song song chưa được phân chia thành các tầng pipeline; khi tăng tần số, đường tổ hợp dài có thể trở thành đường tới hạn và cần được tối ưu theo quy trình timing closure \cite{amd_ug949}. Thứ hai, số liệu công suất hiện dựa trên báo cáo ước tính sau hiện thực, còn tính toàn vẹn tín hiệu mới được đánh giá bằng mô phỏng. Các kết quả này cần được đối chiếu trên phần cứng bằng thiết bị đo phù hợp.